\documentclass[aspectratio=169,11pt]{beamer}
\usepackage{paperdeck}
\title[Attention Is All You Need]{Attention Is All You Need}
\subtitle{Vaswani et al.\ (NeurIPS 2017), arXiv:1706.03762}
\author[Paper summary]{Paper summary for the CS-781 capstone}
\date{}
\begin{document}

\begin{frame}\titlepage\end{frame}

\paperglance
{A.\ Vaswani, N.\ Shazeer, N.\ Parmar, J.\ Uszkoreit, L.\ Jones, A.\ N.\ Gomez, {\L}.\ Kaiser, I.\ Polosukhin. \emph{Attention Is All You Need}. NeurIPS 2017. Google Brain and Google Research.}
{Replace recurrence and convolution entirely with attention: the \textbf{Transformer}, an encoder--decoder built from multi-head self-attention and feed-forward layers.}
{All four of our Phase-1 architectures are Transformers. A2 is closest to this paper's original encoder--decoder design.}
{Main deck: ``The four architectures''. Glossary: ``Transformer building blocks'' and ``Architectures''.}

\begin{frame}{Roadmap}
  \begin{multicols}{2}\small\tableofcontents\end{multicols}
\end{frame}

% ------------------------------------------------------------------
\section{The problem}

\begin{frame}{Why move beyond recurrent networks?}
  \begin{columns}[T]
    \begin{column}{0.52\textwidth}
      \small
      In 2017 the best sequence models were RNNs (LSTMs, GRUs) or CNNs, often joined by attention.
      \vspace{1mm}

      \textbf{RNNs are sequential}: hidden state $h_t$ needs $h_{t-1}$. Positions in one example cannot be processed in parallel. This hurts most on long sequences.
      \vspace{1mm}

      \textbf{CNNs are parallel}, but relating two distant positions needs many layers: $O(n)$ for ConvS2S, $O(\log n)$ for ByteNet.
      \vspace{1mm}

      \textbf{Idea}: attention already relates any two positions in one step. What if we use \emph{only} attention?
    \end{column}
    \begin{column}{0.44\textwidth}
      \begin{exampleblock}{Easy analogy}
        \small
        An RNN reads a 40-page judgment like a clerk reading aloud, one word at a time, remembering everything in one notebook.
        \vspace{1mm}

        A Transformer spreads all 40 pages on a table. Every sentence can glance at every other sentence at once, and many clerks can work in parallel.
      \end{exampleblock}
    \end{column}
  \end{columns}
\end{frame}

\begin{frame}{Self-attention vs recurrence vs convolution}
  \small
  $n$ = sequence length, $d$ = representation size, $k$ = kernel size, $r$ = neighbourhood size.
  \vspace{2mm}

  \begin{tabular}{@{}lccc@{}}
    \toprule
    Layer type & Cost per layer & Sequential operations & Max path length\\
    \midrule
    Self-attention & $O(n^2\cdot d)$ & $O(1)$ & $O(1)$\\
    Recurrent & $O(n\cdot d^2)$ & $O(n)$ & $O(n)$\\
    Convolutional & $O(k\cdot n\cdot d^2)$ & $O(1)$ & $O(\log_k n)$\\
    Restricted self-attention & $O(r\cdot n\cdot d)$ & $O(1)$ & $O(n/r)$\\
    \bottomrule
  \end{tabular}
  \vspace{3mm}

  \begin{itemize}
    \item \textbf{Path length}: how many steps a signal travels between two positions. Shorter paths make long-range dependencies easier to learn.
    \item Self-attention is cheaper than recurrence when $n<d$, which was typical for sentences (e.g.\ $n\approx50$, $d=512$).
    \item For very long inputs, the $n^2$ term dominates. This is why our A4 uses sliding-window (restricted) attention.
  \end{itemize}
\end{frame}

% ------------------------------------------------------------------
\section{The basic idea: attention only}

\begin{frame}{Scaled dot-product attention}
  \begin{columns}[T]
    \begin{column}{0.52\textwidth}
      \small
      Each position has a \textbf{query} $q$, a \textbf{key} $k$ and a \textbf{value} $v$. The output for a query is a weighted sum of the values, weighted by query--key similarity:
      \[
        \text{Attention}(Q,K,V)=\text{softmax}\!\left(\frac{QK^\top}{\sqrt{d_k}}\right)V
      \]
      \textbf{Why divide by $\sqrt{d_k}$?} If $q$ and $k$ have independent unit-variance entries, $q\cdot k$ has variance $d_k$. Large dot products push the softmax into regions with tiny gradients. Scaling keeps the variance at 1.
      \vspace{1mm}

      Dot-product attention is fast because it is just matrix multiplication.
    \end{column}
    \begin{column}{0.44\textwidth}
      \begin{exampleblock}{Easy worked example ($d_k=2$)}
        \small
        Sentence: ``The court dismissed''. Query from ``dismissed'': $q=[1,1]$.
        \vspace{1mm}

        \begin{tabular}{@{}lccc@{}}
          & key & $q\!\cdot\!k/\sqrt2$ & weight\\
          \midrule
          The & $[0,1]$ & 0.71 & 0.16\\
          court & $[1,2]$ & 2.12 & \textbf{0.67}\\
          dismissed & $[1,0]$ & 0.71 & 0.16\\
        \end{tabular}
        \vspace{1mm}

        With values The $[1,0]$, court $[0,1]$, dismissed $[1,1]$, the output is $0.16[1,0]+0.67[0,1]+0.16[1,1]=[0.33,\,0.84]$.
        ``Dismissed'' mostly gathers information from ``court'', \emph{who} dismissed.
      \end{exampleblock}
    \end{column}
  \end{columns}
\end{frame}

\begin{frame}{Multi-head attention}
  \begin{columns}[T]
    \begin{column}{0.52\textwidth}
      \small
      Instead of one attention over $d_{\text{model}}=512$ dimensions, run $h=8$ \textbf{heads} in parallel, each on its own learned projections of size $d_k=d_v=64$:
      \[
        \text{head}_i=\text{Attention}(QW_i^Q,\;KW_i^K,\;VW_i^V)
      \]
      \[
        \text{MultiHead}=\text{Concat}(\text{head}_1,\ldots,\text{head}_h)\,W^O
      \]
      Total cost is about that of one full-size head.
      \vspace{1mm}

      \textbf{Why}: one head averages everything into one weighted sum. Several heads can attend to \emph{different} things at once.
    \end{column}
    \begin{column}{0.44\textwidth}
      \begin{exampleblock}{Intuition in a judgment}
        \small
        In ``The respondent, who was represented by counsel, did not appear before the tribunal'':
        \begin{itemize}
          \item one head links ``did not appear'' to ``respondent'' (subject);
          \item another links ``appear'' to ``tribunal'' (location);
          \item another links ``represented'' to ``counsel''.
        \end{itemize}
        The paper's appendix shows heads that track syntax and coreference.
      \end{exampleblock}
    \end{column}
  \end{columns}
\end{frame}

\begin{frame}{Three uses of attention in the model}
  \begin{columns}[T]
    \begin{column}{0.5\textwidth}
      \small
      \begin{enumerate}
        \item \textbf{Encoder self-attention}: every source position attends to every source position.
        \item \textbf{Decoder masked self-attention}: each target position attends only to \emph{earlier} target positions. Future scores are set to $-\infty$ before the softmax, so generation stays auto-regressive.
        \item \textbf{Encoder--decoder (cross) attention}: decoder queries attend to the encoder's outputs, the keys and values. Each output word can look at the whole input.
      \end{enumerate}
    \end{column}
    \begin{column}{0.46\textwidth}
      \begin{exampleblock}{Translation example}
        \small
        Translating ``The appeal is dismissed'' into Hindi:
        \begin{itemize}
          \item encoder self-attention works out that ``dismissed'' refers to ``appeal'';
          \item while writing each Hindi word, the decoder's masked self-attention sees the Hindi words already written;
          \item cross-attention looks back at the English words it needs next.
        \end{itemize}
      \end{exampleblock}
    \end{column}
  \end{columns}
\end{frame}

\begin{frame}{The full architecture}
  \begin{columns}[T]
    \begin{column}{0.46\textwidth}
      \centering
\resizebox{0.95\linewidth}{!}{%
      \begin{tikzpicture}[node distance=2mm]
        \node[box,minimum width=28mm] (ein) {input emb.\ + positions};
        \node[hot,above=of ein,minimum width=28mm] (esa) {multi-head self-attn};
        \node[box,above=of esa,minimum width=28mm] (eff) {feed-forward};
        \node[draw=navy,dashed,rounded corners,fit=(esa)(eff),inner sep=3pt,label={[font=\scriptsize]left:$\times6$}] (E) {};
        \node[box,right=9mm of ein,minimum width=28mm] (din) {output emb.\ + positions};
        \node[hot,above=of din,minimum width=28mm] (dsa) {masked self-attn};
        \node[hot,above=of dsa,minimum width=28mm] (dca) {cross-attn};
        \node[box,above=of dca,minimum width=28mm] (dff) {feed-forward};
        \node[draw=navy,dashed,rounded corners,fit=(dsa)(dff),inner sep=3pt,label={[font=\scriptsize]right:$\times6$}] (D) {};
        \node[box,above=4mm of dff] (out) {linear + softmax};
        \draw[arr] (ein)--(esa); \draw[arr] (din)--(dsa); \draw[arr] (dff)--(out);
        \draw[arr,warm] (eff.east) -- node[above,font=\tiny]{memory} (dca.west);
      \end{tikzpicture}}
    \end{column}
    \begin{column}{0.5\textwidth}
      \small
      \begin{itemize}
        \item Each sub-layer is wrapped as $\text{LayerNorm}(x+\text{Sublayer}(x))$: a residual connection, then layer normalisation (\emph{post-LN}).
        \item Feed-forward network, applied to each position separately:
          \[
            \text{FFN}(x)=\max(0,xW_1+b_1)W_2+b_2
          \]
          with $d_{ff}=2048$ (4$\times$ wider).
        \item Input and output embeddings and the pre-softmax layer share one weight matrix; embeddings are multiplied by $\sqrt{d_{\text{model}}}$.
        \item Base model: $N=6$, $d_{\text{model}}=512$, $h=8$: 65M parameters.
      \end{itemize}
    \end{column}
  \end{columns}
\end{frame}

\begin{frame}{Positional encoding}
  \begin{columns}[T]
    \begin{column}{0.52\textwidth}
      \small
      Attention ignores order, so position information is \emph{added} to the embeddings:
      \[ PE_{(pos,2i)}=\sin\!\big(pos/10000^{2i/d}\big) \]
      \[ PE_{(pos,2i+1)}=\cos\!\big(pos/10000^{2i/d}\big) \]
      \begin{itemize}
        \item Wavelengths form a geometric series from $2\pi$ to $10000\cdot2\pi$.
        \item For any fixed offset $k$, $PE_{pos+k}$ is a \emph{linear function} of $PE_{pos}$. The authors hoped this would let the model attend by relative position.
        \item Learned position embeddings gave nearly identical results. Sinusoids were chosen because they might extrapolate to longer sequences.
      \end{itemize}
    \end{column}
    \begin{column}{0.44\textwidth}
      \begin{exampleblock}{Easy analogy: a clock}
        \small
        Each pair of dimensions is a clock hand turning at its own speed. Fast hands tell neighbouring positions apart; slow hands tell distant ones apart. Together they give each position a unique reading.
      \end{exampleblock}
      \begin{block}{In our models}
        \small
        A2 uses these sinusoids. A1 uses learned positions. A3 (ALiBi) and A4 (RoPE) use later relative schemes.
      \end{block}
    \end{column}
  \end{columns}
\end{frame}

\begin{frame}{Training recipe}
  \begin{columns}[T]
    \begin{column}{0.52\textwidth}
      \small
      \begin{itemize}
        \item \textbf{Data}: WMT 2014 EN--DE (4.5M sentence pairs, shared BPE vocabulary of about 37k); EN--FR (36M sentences, 32k word-pieces). Batches of about 25k source and 25k target tokens.
        \item \textbf{Hardware}: 8 NVIDIA P100 GPUs. Base: 100k steps in 12 hours. Big: 300k steps in 3.5 days.
        \item \textbf{Adam} with $\beta_1=0.9$, $\beta_2=0.98$, $\epsilon=10^{-9}$.
        \item \textbf{Learning rate} with warm-up, then inverse-square-root decay:
        \[
          \text{lr}=d_{\text{model}}^{-0.5}\min\!\big(\text{step}^{-0.5},\,\text{step}\cdot\text{warmup}^{-1.5}\big)
        \]
        with warmup $=4000$.
      \end{itemize}
    \end{column}
    \begin{column}{0.44\textwidth}
      \begin{block}{Regularisation}
        \small
        \begin{itemize}
          \item Dropout 0.1 on each sub-layer output and on embedding sums.
          \item \textbf{Label smoothing} $\epsilon=0.1$: the target puts 0.9 on the true token and spreads 0.1 over the rest. This hurts perplexity but improves BLEU.
        \end{itemize}
      \end{block}
      \begin{exampleblock}{Inference}
        \small
        Beam search (beam 4, length penalty 0.6). Averages of the last 5 (base) or 20 (big) checkpoints.
      \end{exampleblock}
    \end{column}
  \end{columns}
\end{frame}

% ------------------------------------------------------------------
\section{Results}

\begin{frame}{Machine translation (newstest2014)}
  \begin{columns}[T]
    \begin{column}{0.56\textwidth}
      \footnotesize
      \begin{tabular}{@{}lcccc@{}}
        \toprule
        & \multicolumn{2}{c}{BLEU} & \multicolumn{2}{c}{Training FLOPs}\\
        Model & EN-DE & EN-FR & EN-DE & EN-FR\\
        \midrule
        GNMT + RL & 24.6 & 39.92 & $2.3\!\cdot\!10^{19}$ & $1.4\!\cdot\!10^{20}$\\
        ConvS2S & 25.16 & 40.46 & $9.6\!\cdot\!10^{18}$ & $1.5\!\cdot\!10^{20}$\\
        MoE & 26.03 & 40.56 & $2.0\!\cdot\!10^{19}$ & $1.2\!\cdot\!10^{20}$\\
        ConvS2S ensemble & 26.36 & 41.29 & $7.7\!\cdot\!10^{19}$ & $1.2\!\cdot\!10^{21}$\\
        \midrule
        Transformer (base) & 27.3 & 38.1 & \multicolumn{2}{c}{$3.3\!\cdot\!10^{18}$}\\
        Transformer (big) & \textbf{28.4} & \textbf{41.8} & \multicolumn{2}{c}{$2.3\!\cdot\!10^{19}$}\\
        \bottomrule
      \end{tabular}
    \end{column}
    \begin{column}{0.4\textwidth}
      \begin{exampleblock}{Reading the table}
        \small
        \begin{itemize}
          \item EN--DE: 28.4 BLEU, \textbf{more than 2 BLEU} above the best earlier models, including ensembles.
          \item EN--FR: a new single-model record of 41.8.
          \item The base model already beats all earlier EN--DE systems with a fraction of the training compute.
        \end{itemize}
      \end{exampleblock}
      {\scriptsize The paper's text states 41.0 for EN--FR in Section 6.1; the abstract and Table 2 give 41.8.}
    \end{column}
  \end{columns}
\end{frame}

\begin{frame}{Model variations (EN--DE dev set, newstest2013)}
  \begin{columns}[T]
    \begin{column}{0.5\textwidth}
      \footnotesize
      \begin{tabular}{@{}lccc@{}}
        \toprule
        Change from base & PPL & BLEU & Params\\
        \midrule
        Base ($h=8$, $d_k=64$) & 4.92 & 25.8 & 65M\\
        $h=1$, $d_k=512$ & 5.29 & 24.9 & \\
        $h=4$, $d_k=128$ & 5.00 & 25.5 & \\
        $h=16$, $d_k=32$ & 4.91 & 25.8 & \\
        $h=32$, $d_k=16$ & 5.01 & 25.4 & \\
        $d_k=16$ (smaller keys) & 5.16 & 25.1 & 58M\\
        $N=2$ layers & 6.11 & 23.7 & 36M\\
        dropout 0.0 & 5.77 & 24.6 & \\
        learned positions & 4.92 & 25.7 & \\
        Big ($d=1024$, $h=16$) & \textbf{4.33} & \textbf{26.4} & 213M\\
        \bottomrule
      \end{tabular}
    \end{column}
    \begin{column}{0.46\textwidth}
      \begin{exampleblock}{Lessons}
        \small
        \begin{itemize}
          \item One head is 0.9 BLEU worse; too many heads also hurt.
          \item Smaller keys hurt: judging compatibility is not trivial.
          \item Bigger models are better.
          \item Dropout matters a lot for avoiding overfitting.
          \item Learned and sinusoidal positions are about equal.
        \end{itemize}
      \end{exampleblock}
      {\scriptsize Perplexities are per word-piece.}
    \end{column}
  \end{columns}
\end{frame}

\begin{frame}{Beyond translation: constituency parsing}
  \begin{columns}[T]
    \begin{column}{0.52\textwidth}
      \small
      Parsing turns a sentence into a syntax tree, written as a longer bracketed sequence. It has strong structural constraints, and RNN sequence-to-sequence models struggled with small data.
      \vspace{1mm}

      A 4-layer Transformer ($d_{\text{model}}=1024$), with almost no task-specific tuning, on WSJ Section 23:
      \vspace{1mm}

      \begin{tabular}{@{}lc@{}}
        \toprule
        Setting & F1\\
        \midrule
        WSJ only (40K sentences) & 91.3\\
        Semi-supervised (about 17M sentences) & 92.7\\
        \bottomrule
      \end{tabular}
    \end{column}
    \begin{column}{0.44\textwidth}
      \begin{exampleblock}{Why it matters}
        \small
        With only the 40K WSJ sentences, the Transformer beat the Berkeley parser (90.4), where RNN sequence-to-sequence models had failed.
        Only one model in the comparison, the RNN Grammar, did better.
        \vspace{1mm}

        This was early evidence that the architecture is general-purpose, as our seven phases assume.
      \end{exampleblock}
    \end{column}
  \end{columns}
\end{frame}

% ------------------------------------------------------------------
\section{Assessment and outlook}

\begin{frame}{How our Phase-1 models descend from this paper}
  \small
  \begin{tabularx}{\textwidth}{@{}lXX@{}}
    \toprule
    Paper component & Our A1--A4 & Change and reason\\
    \midrule
    Encoder--decoder & A2 keeps it; A1, A4 decoder-only; A3 one shared stack & decoder-only is simpler for language modelling\\
    Post-LN & pre-LN (A1--A3), RMSNorm (A4) & pre-LN trains more stably when deep\\
    ReLU FFN & ReLU (A2), GELU (A1, A3), SwiGLU (A4) & gated units work better (Shazeer, 2020)\\
    Sinusoidal positions & sinusoids (A2), learned (A1), ALiBi (A3), RoPE (A4) & relative schemes extrapolate better\\
    Multi-head attention & MHA; grouped-query in A4 & smaller KV-cache\\
    Warm-up + inverse sqrt & warm-up + cosine & common modern choice\\
    Tied embeddings & tied embeddings & fewer parameters\\
    \bottomrule
  \end{tabularx}
\end{frame}

\begin{frame}{Strengths and limitations}
  \begin{columns}[T]
    \begin{column}{0.48\textwidth}
      \begin{exampleblock}{Strengths}
        \small
        \begin{itemize}
          \item Fully parallel training over positions.
          \item Constant path length between any two tokens.
          \item Better quality at a fraction of the training cost.
          \item Simple, uniform building blocks; scales with size.
          \item Attention weights are partly interpretable.
        \end{itemize}
      \end{exampleblock}
    \end{column}
    \begin{column}{0.48\textwidth}
      \begin{alertblock}{Limitations}
        \small
        \begin{itemize}
          \item $O(n^2)$ time and memory in sequence length. A problem for long legal judgments.
          \item Generation is still one token at a time.
          \item Positions only via added encodings; extrapolation to longer inputs was not tested.
          \item Post-LN can be unstable in deep stacks without warm-up.
          \item Evaluated mainly on translation and one parsing task.
        \end{itemize}
      \end{alertblock}
    \end{column}
  \end{columns}
\end{frame}

\begin{frame}{Future research directions (1): stated in the paper, and what followed}
  \small
  \begin{tabularx}{\textwidth}{@{}lX@{}}
    \toprule
    Direction & What followed\\
    \midrule
    Other modalities \emph{(paper)} & Vision Transformer (Dosovitskiy et al., 2021) for images; Transformers for speech and video\\
    Local, restricted attention for large inputs \emph{(paper)} & Longformer (Beltagy et al., 2020) and sliding windows as in Mistral; FlashAttention (Dao et al., 2022) makes exact attention memory-efficient\\
    Less sequential generation \emph{(paper)} & Non-autoregressive translation; speculative decoding\\
    Pre-training at scale & BERT (encoder-only), GPT (decoder-only), T5 (encoder--decoder)\\
    Better positions & Relative positions (Shaw et al., 2018), RoPE (2021), ALiBi (2022)\\
    Normalisation and FFN & Pre-LN (Xiong et al., 2020); gated FFNs such as SwiGLU (Shazeer, 2020)\\
    \bottomrule
  \end{tabularx}
\end{frame}

\begin{frame}{Future research directions (2): open questions}
  \begin{columns}[T]
    \begin{column}{0.48\textwidth}
      \begin{block}{Open in general}
        \small
        \begin{itemize}
          \item Sub-quadratic sequence models that match attention quality, such as state-space models.
          \item Why attention heads specialise, and which heads matter.
          \item Scaling laws: the best balance of model size, data and compute.
          \item Length generalisation beyond training context.
        \end{itemize}
      \end{block}
    \end{column}
    \begin{column}{0.48\textwidth}
      \begin{exampleblock}{Open for legal NLP}
        \small
        \begin{itemize}
          \item Reading whole judgments (tens of thousands of tokens) within a free GPU's memory.
          \item Does the encoder--decoder split (A2) help summarisation more than decoder-only models?
          \item Can attention maps serve as explanations a lawyer would accept?
        \end{itemize}
      \end{exampleblock}
    \end{column}
  \end{columns}
\end{frame}

\begin{frame}{Summary}
  \begin{itemize}
    \item The Transformer removes recurrence and convolution. It uses \textbf{scaled dot-product attention}, \textbf{multi-head} attention, position-wise feed-forward layers, residual connections and LayerNorm.
    \item Masked decoder self-attention keeps generation auto-regressive; cross-attention links decoder and encoder.
    \item Sinusoidal positional encodings supply word order.
    \item Results: 28.4 BLEU (EN--DE) and 41.8 (EN--FR), state of the art at much lower training cost. It also generalises to parsing.
    \item Every model in our capstone is a descendant of this architecture.
  \end{itemize}
\end{frame}

\begin{frame}{Reference}
  \small
  Vaswani, A., Shazeer, N., Parmar, N., Uszkoreit, J., Jones, L., Gomez, A.\ N., Kaiser, {\L}., and Polosukhin, I.\ (2017). Attention Is All You Need. In \emph{Advances in Neural Information Processing Systems 30}. arXiv:1706.03762.
  \vspace{3mm}

  {\footnotesize Later work cited: Shaw et al.\ (2018); Beltagy et al.\ (2020) Longformer; Xiong et al.\ (2020) On Layer Normalization in the Transformer Architecture; Shazeer (2020) GLU Variants; Dosovitskiy et al.\ (2021) ViT; Dao et al.\ (2022) FlashAttention.}
\end{frame}

\end{document}
